% =====================================================================
%  Bölüm 13 — Ardışık Hesap ve Kesme Eleme
% =====================================================================
\chapter{Ardışık Hesap ve Kesme Eleme}\label{bol:ardisik}

\begin{bolumhedefleri}
\begin{itemize}
  \item Gentzen'in ardışık hesabını, LK ve LJ sistemlerini ve aralarındaki
        farkı,
  \item Çok sonuçlu ardışıkların klasik mantığı nasıl ``kendiliğinden''
        verdiğini,
  \item Kuralların tersinirliğini ve bunun bir ispat arama algoritmasına
        nasıl dönüştüğünü,
  \item Kesme kuralını, Gentzen'in \emph{Hauptsatz}'ını ve kesmenin
        \Lean{}'de anlambilimsel yoldan elenmesini,
  \item Alt formül özelliğini ve bunun tutarlılık, karar verilebilirlik ve
        otomatik ispat için sonuçlarını.
\end{itemize}
\end{bolumhedefleri}

Kısım~III'ün ilk üç bölümünde önermeler mantığının \emph{bir} ispat
sistemini, doğal çıkarımı, inceledik ve onun iki değerli anlambilime göre
hem geçerli hem tam olduğunu gösterdik. Bu bölümde aynı mantık için
ikinci bir ispat sistemi tanıtıyoruz: Gerhard Gentzen'in \emph{ardışık
hesabı}. İki sistem aynı formülleri türetir. Ama ardışık hesabın
kurallarının biçimi ona doğal çıkarımda olmayan bir özellik kazandırır:
Türetimler \emph{aşağıdan yukarıya}, mekanik olarak aranabilir. Bu bölümün
sonunda, \Lean{}'de çalışan küçük bir otomatik ispatlayıcımız olacak.

% ---------------------------------------------------------------------
\section{Ardışıklar}\label{sec:ardisiklar}

Bölüm~5'te doğal çıkarım kurallarını $\Gamma \vdash A$ biçimindeki
\emph{ardışıklarla} yazmıştık. Solda bir varsayım listesi, sağda tek bir
sonuç vardı. Gentzen'in ardışık hesabında sağ taraf da bir listedir:

\begin{tanim}{Ardışık}{ardisik}
Bir \textbf{ardışık}, iki formül listesinden oluşan bir çifttir ve
\[
  A_1, \dots, A_m \;\Rightarrow\; B_1, \dots, B_n
\]
diye yazılır. Anlamı şudur: \emph{Soldaki formüllerin hepsi doğruysa,
sağdaki formüllerden en az biri doğrudur.} Yani sol taraf bir ``ve'', sağ
taraf bir ``veya'' gibi okunur:
\[
  \Gamma \Rightarrow \Delta
  \quad\text{geçerlidir}\quad\iff\quad
  \models \textstyle\bigwedge\Gamma \oise \bigvee\Delta .
\]
\end{tanim}

\noindent İki uç durum öğreticidir. $\Rightarrow A$ (sol boş), $A$'nın bir
totoloji olduğunu söyler. $A \Rightarrow$ (sağ boş) ise $A$'nın
\emph{sağlanamaz} olduğunu söyler, çünkü boş bir ``veya'' hiçbir zaman
doğru değildir. Boş ardışık $\Rightarrow$ ise hiçbir zaman geçerli
olamaz; onu türetebilen bir sistem tutarsızdır.

\begin{tarihce}[Tarihçe: Gentzen ve Hauptsatz]
Gerhard Gentzen, 1934--35'te yayımlanan ``Mantıksal çıkarım üzerine
incelemeler'' adlı tezinde \citep{gentzen1935} iki sistem birden tanıttı:
\emph{doğal çıkarım} (NK, NJ) ve \emph{ardışık hesap} (LK, LJ). Doğal
çıkarımın matematikçilerin gerçek akıl yürütmesine yakın olduğunu, ama
onun hakkında teoremler ispatlamanın zor olduğunu gördü. Ardışık hesabı bu
yüzden geliştirdi: Tezin ana sonucu olan \emph{Hauptsatz} (ana teorem),
ardışık hesapta \emph{kesme} kuralının elenebileceğini söyler. Gentzen bir
yıl sonra aynı teknikleri kullanarak Peano aritmetiğinin tutarlılığını
ispatladı \citep{gentzen1936}. Ardışık hesabın bugünkü standart sunumu
için \citet{troelstra2000} ve \citet{negri2001} başvuru kaynaklarıdır.
\end{tarihce}

% ---------------------------------------------------------------------
\section{Kesmesiz bir ardışık hesap}\label{sec:g3}

Ardışık hesabın her bağlaç için iki kuralı vardır: bağlacı \emph{solda}
tanıtan bir kural ve \emph{sağda} tanıtan bir kural. Eleme kuralı yoktur.
Bunun yerine doğal çıkarımdaki eleme kurallarının işini sol kurallar
görür. \Cref{fig:g3-kurallar}, \Lean{}'de tanımlayacağımız sistemin
kurallarını gösteriyor.

\begin{figure}[htbp]
\centering
\setlength{\tabcolsep}{14pt}
\renewcommand{\arraystretch}{3.0}
\begin{tabular}{cc}
\begin{prooftree}
  \hypo{p_n \in \Gamma}
  \hypo{p_n \in \Delta}
  \infer2[Aks]{\Gamma \Rightarrow \Delta}
\end{prooftree}
&
\begin{prooftree}
  \hypo{\bot \in \Gamma}
  \infer1[$\bot$L]{\Gamma \Rightarrow \Delta}
\end{prooftree}
\\
\begin{prooftree}
  \hypo{A, B, \Gamma \Rightarrow \Delta}
  \infer1[$\wedge$L]{A \ove B, \Gamma \Rightarrow \Delta}
\end{prooftree}
&
\begin{prooftree}
  \hypo{\Gamma \Rightarrow A, \Delta}
  \hypo{\Gamma \Rightarrow B, \Delta}
  \infer2[$\wedge$R]{\Gamma \Rightarrow A \ove B, \Delta}
\end{prooftree}
\\
\begin{prooftree}
  \hypo{A, \Gamma \Rightarrow \Delta}
  \hypo{B, \Gamma \Rightarrow \Delta}
  \infer2[$\vee$L]{A \oveya B, \Gamma \Rightarrow \Delta}
\end{prooftree}
&
\begin{prooftree}
  \hypo{\Gamma \Rightarrow A, B, \Delta}
  \infer1[$\vee$R]{\Gamma \Rightarrow A \oveya B, \Delta}
\end{prooftree}
\\
\begin{prooftree}
  \hypo{\Gamma \Rightarrow A, \Delta}
  \hypo{B, \Gamma \Rightarrow \Delta}
  \infer2[$\to$L]{A \oise B, \Gamma \Rightarrow \Delta}
\end{prooftree}
&
\begin{prooftree}
  \hypo{A, \Gamma \Rightarrow B, \Delta}
  \infer1[$\to$R]{\Gamma \Rightarrow A \oise B, \Delta}
\end{prooftree}
\\
\multicolumn{2}{c}{%
\colorbox{fikirbg}{\begin{prooftree}
  \hypo{\Gamma \Rightarrow \Delta}
  \hypo{\Gamma \subseteq \Gamma'}
  \hypo{\Delta \subseteq \Delta'}
  \infer3[Yapı]{\Gamma' \Rightarrow \Delta'}
\end{prooftree}\quad\raisebox{0.4em}{\textcolor{fikirframe}{\baslikfont\scriptsize\bfseries yapısal kural}}}}
\\
\end{tabular}
\caption{Kesmesiz klasik ardışık hesap. Her bağlaç için bir sol (L) ve
bir sağ (R) kuralı vardır. Son kural, Gentzen'in zayıflatma, değiştirme
ve daraltma kurallarını tek bir kuralda toplar. Kesme kuralı \emph{yoktur}.}
\label{fig:g3-kurallar}
\end{figure}

\noindent Kurallara yakından bakalım. $\to$L kuralı, doğal çıkarımdaki
modus ponensin ardışık hesaptaki karşılığıdır. Solda $A \oise B$ varsa ve
$A$'yı ``sağlayabiliyorsak'' (sol öncül), $B$'yi varsayabiliriz (sağ öncül).
$\vee$R kuralı ise çok sonuçlu olmanın gücünü gösterir: $A \oveya B$'yi
kanıtlamak için hangisinin doğru olduğuna \emph{karar vermek} zorunda
değiliz; ikisini birden sağa yazarız. \Lean{}'deki tanım şöyle:

\leandosya[Kesmesiz ardışık hesap]{Bolum13}{ardisik}

\noindent \lc{Ardisik Γ Δ} tanımı, Bölüm~11'deki \lc{Turetilir}
tanımıyla aynı kalıptadır: Her kural bir kurucudur. Tek fark, ana
formülün listenin \emph{başında} durmasıdır. Ana formül başka bir yerdeyse,
\lc{degis} yapısal kuralıyla öne alınır.

\begin{dikkat}[Dikkat: \lc{::} ve bağlaçların önceliği]
\Lean{}'de \lc{::} işlecinin önceliği (67), bağlaçlarımızınkinden
(25--35) yüksektir. Bu yüzden \lc{A {\sembolkodfont ⋀} B :: Γ}, \lc{A {\sembolkodfont ⋀} (B :: Γ)} diye
okunur ve bir tip hatası verir. Ana formülü her zaman paranteze alıyoruz:
\lc{(A {\sembolkodfont ⋀} B) :: Γ}.
\end{dikkat}

\subsection{LK ve LJ}

Gentzen'in klasik sistemi LK'dır. Sezgici sistem LJ ise LK'dan tek bir
kısıtlamayla elde edilir: \emph{Sağ tarafta en fazla bir formül
bulunabilir.} Bu küçük sözdizimsel kısıtlama, Bölüm~4'te uzun uzun
tartıştığımız yapıcı/klasik ayrımını tam olarak yakalar. Bunu görmek için
Peirce yasasının türetimine bakalım (\cref{fig:peirce-ardisik}).

\leandosya[İki örnek türetim]{Bolum13}{ornekler}

\begin{figure}[htbp]
\centering
\begin{prooftree}
  \infer0[Aks]{p \Rightarrow \colorbox{teorembg}{$q,\; p$}}
  \infer1[$\to$R]{\Rightarrow \colorbox{teorembg}{$p \oise q,\; p$}}
  \infer0[Aks]{p \Rightarrow p}
  \infer2[$\to$L]{(p \oise q) \oise p \Rightarrow p}
  \infer1[$\to$R]{\Rightarrow ((p \oise q) \oise p) \oise p}
\end{prooftree}
\caption{Peirce yasasının ardışık hesaptaki türetimi, \Lean{}'deki
\lc{peirce\_ardisik} ispatının adım adım karşılığı. Renkli ardışıklarda
sağ tarafta \emph{iki} formül vardır; LJ'de bu adımlara izin verilmez. Bu
yüzden Peirce yasası LJ'de türetilemez.}
\label{fig:peirce-ardisik}
\end{figure}

\begin{onemlifikir}[Önemli Fikir: Klasik mantık = çok sonuçlu ardışıklar]
Doğal çıkarımda klasik mantığa geçmek için yeni bir \emph{kural}
eklememiz gerekmişti (RAA). Ardışık hesapta kurallar aynı kalır; değişen
yalnızca sağ tarafın kaç formül taşıyabildiğidir. Klasik mantığın ``ya
öyle ya böyle, ama hangisi olduğunu söylemeden'' akıl yürütmesi, sağ
taraftaki virgülde saklıdır. Girard'ın deyişiyle, klasik ve sezgici mantık
arasındaki fark bağlaçlarda değil, \emph{yapıdadır} \citep{girard1989}.
\end{onemlifikir}

% ---------------------------------------------------------------------
\section{Anlambilim ve geçerlilik}\label{sec:ardisik-gecerlilik}

Bir ardışığın anlamını \Lean{}'de iki küçük yüklemle ifade ediyoruz. Sol
tarafın hepsi doğru (\lc{HepsiDogru}), sağ taraftan biri doğru
(\lc{BiriDogru}):

\leandosya[Ardışıkların anlamı]{Bolum13}{gecerli_tanim}

\noindent Listelerin başına bir eleman eklendiğinde bu iki yüklemin nasıl
davrandığını söyleyen iki önsav, bölümün geri kalanında sürekli
kullanılacak. Üçüncü önsav ise geçerliliğin listeler büyüdükçe
korunduğunu söyler (bu, yapısal kuralın geçerliliğidir):

\leandosya[Liste önsavları]{Bolum13}{liste_onsavlari}

\noindent Geçerlilik teoremi Bölüm~11'deki gibi kural tümevarımıyla
ispatlanır. Her kural, öncüllerindeki geçerliliği sonucuna taşır:

\leandosya[Ardışık hesabın geçerliliği]{Bolum13}{ardisik_gecerli}

\noindent Son satırdaki \lc{bos\_ardisik\_yok}, boş ardışığın
türetilemediğini, yani sistemin tutarlı olduğunu söyler. Geçerlilik
ispatı burada da kısa bir sonuç verir: Boş ardışık hiçbir değerlemede
doğru değildir.

% ---------------------------------------------------------------------
\section{Tersinirlik ve ispat arama}\label{sec:tersinirlik}

Doğal çıkarımda bir formülü türetmeye çalışırken çoğu zaman ne yapacağımızı
\emph{tahmin etmemiz} gerekir. Modus ponens hangi $A$ ile uygulanmalı?
Hangi ara önsav işe yarar? Ardışık hesabın kurallarında böyle bir tahmin
yoktur. Her kuralın öncüllerinde yalnızca sonucunda geçen formüllerin
\emph{parçaları} bulunur. Üstelik bizim sistemimizin her kuralı
\emph{tersinirdir}:

\begin{tanim}{Tersinir kural}{tersinir}
Bir kural, sonucu geçerli olduğunda öncüllerinin hepsi de geçerliyse
\textbf{tersinirdir}.
\end{tanim}

\noindent Tersinirlik, ispat aramada hiçbir zaman yanlış bir seçim
yapamayacağımız anlamına gelir. Geçerli bir ardışığa \emph{hangi} kuralı
uygularsak uygulayalım, öncüller de geçerli kalır. Bu yüzden arama hiçbir
zaman geri dönmek (\emph{backtracking}) zorunda kalmaz.

\leandosya[Tersinirlik önsavları]{Bolum13}{tersinirlik}

\begin{figure}[htbp]
\centering
\begin{tikzpicture}[
    ard/.style={draw=#1, fill=#1!8, thick, rounded corners=3pt,
                font=\small, inner sep=4pt},
    ok/.style={-{Stealth[length=2.2mm]}, thick, black!45},
    kural/.style={font=\baslikfont\scriptsize\bfseries, text=leantac},
  ]
  \node[ard=teoremframe] (k) at (0,0) {$\Rightarrow (p \oveya q) \oise (q \oveya p)$};
  \node[ard=anarenk] (a1) at (0,1.4) {$p \oveya q \Rightarrow q \oveya p$};
  \node[ard=anarenk] (a2) at (0,2.8) {$p \oveya q \Rightarrow q,\; p$};
  \node[ard=ikincil] (l) at (-2.6,4.3) {$p \Rightarrow q,\; p$};
  \node[ard=ikincil] (r) at (2.6,4.3) {$q \Rightarrow q,\; p$};
  \node[font=\baslikfont\scriptsize\bfseries, text=ikincil] at (-2.6,5.0) {Aks ($p$)};
  \node[font=\baslikfont\scriptsize\bfseries, text=ikincil] at (2.6,5.0) {Aks ($q$)};
  \draw[ok] (k) -- node[right, kural]{$\to$R} (a1);
  \draw[ok] (a1) -- node[right, kural]{$\vee$R} (a2);
  \draw[ok] (a2) -- node[left, kural, pos=0.45]{$\vee$L} (l);
  \draw[ok] (a2) -- (r);
  \node[font=\baslikfont\scriptsize, text=black!55, align=left, anchor=west] at (4.6,1.4)
       {her adımda:\\en soldaki bağlacı\\parçala};
  \node[font=\baslikfont\scriptsize, text=black!55, align=left, anchor=west] at (4.6,4.3)
       {yapraklar:\\yalnızca atomlar};
  \draw[-{Stealth[length=3mm]}, very thick, black!25] (-5.2,0) -- (-5.2,4.6)
       node[midway, left, font=\baslikfont\scriptsize, text=black!55, align=right]{arama\\yönü};
\end{tikzpicture}
\caption{$(p \oveya q) \oise (q \oveya p)$ için ispat arama. Kök, ispatlamak
istediğimiz ardışıktır. Her adımda bir bağlaç parçalanır. Yapraklarda
yalnızca atomlar kaldığında, solda ve sağda ortak bir atom olup olmadığına
bakılır.}
\label{fig:ispat-arama}
\end{figure}

\subsection{Kesmesiz tamlık}

Tersinirlik, bir tamlık ispatına doğrudan dönüşür. Geçerli bir ardışık
verilsin. Bağlaçları birer birer parçalayalım; tersinirlik sayesinde her
yeni ardışık da geçerlidir. Bağlaç kalmadığında elimizde yalnızca
atomlardan oluşan geçerli ardışıklar kalır. Böyle bir ardışık ya solda ve
sağda ortak bir atom içerir (o zaman Aks kuralı uygulanır), ya da
geçersizdir. Çünkü soldaki atomları doğru, diğer her şeyi yanlış yapan
değerleme bir karşı model olurdu.

Bunu \Lean{}'de yazmak için önce aramanın sona ereceğini garanti eden bir
\emph{ölçü} tanımlıyoruz. Her kural uygulaması bu ölçüyü en az bir
azaltır:

\leandosya[Ardışığın ölçüsü]{Bolum13}{olcu}

\noindent Ölçüdeki $2 \cdot \mathrm{boyut}$ çarpanı önemlidir. Örneğin
$\wedge$L kuralı $A \ove B$'yi $A$ ve $B$ ile değiştirir. Bağlaç sayısı
bir azalır ama formül sayısı bir artar. Bağlaçlara iki kat ağırlık
verilmezse toplam değişmezdi.

Atomik formülleri işlerken onları listenin sonundaki bir ``atom
bölgesine'' taşırız. \lc{tasi1} ve \lc{tasi2}, bu taşımanın üyeliği
değiştirmediğini söyler. Böylece \lc{degis} kuralı taşımayı türetime
dönüştürebilir:

\leandosya[Atomları kenara almak]{Bolum13}{tasi}

\noindent Aramanın yaprakları, yalnızca atomlardan oluşan ardışıklardır.
Burada karşı model yöntemi devreye girer:

\leandosya[Atomik ardışıklar]{Bolum13}{atomik}

\noindent Artık arama algoritmasını bir tamlık ispatı olarak yazabiliriz.
\lc{tam\_aux}, ölçü üzerinde (\lc{n} ile sınırlanmış) özyinelemeyle
tanımlanmış bir teoremdir. Önce soldaki formülleri, sol boşalınca da
sağdakileri işler. Her durumda tersinirlik önsavı yeni ardışığın geçerli
olduğunu, özyinelemeli çağrı ise türetilebildiğini söyler:

\leandosya[İspat arama, bir tamlık ispatı olarak]{Bolum13}{tam_aux}

\noindent \lc{tam\_aux}'un örüntüleri, \cref{fig:ispat-arama}'daki
aramanın birebir tarifidir. \lc{P} ve \lc{Q} kenara alınmış atomları
tutar. Ardından son adım gelir: Başlangıçta kenarda hiç atom yoktur.

\leandosya[Kesmesiz tamlık]{Bolum13}{ardisik_tam}

\begin{teorem}{Kesmesiz tamlık}{kesmesiz-tamlik}
Geçerli her $\Gamma \Rightarrow \Delta$ ardışığı, kesme kuralı
kullanılmadan türetilebilir.
\end{teorem}

% ---------------------------------------------------------------------
\section{Kesme kuralı}\label{sec:kesme}

Matematikte bir teoremi ispatlarken önce bir \emph{önsav} ispatlar, sonra
onu kullanırız. Ardışık hesapta bu adımın adı \emph{kesme}dir:
\begin{center}
\begin{prooftree}
  \hypo{\Gamma \Rightarrow A, \Delta}
  \hypo{A, \Gamma \Rightarrow \Delta}
  \infer2[\textcolor{teoremframe}{Kesme}]{\Gamma \Rightarrow \Delta}
\end{prooftree}
\end{center}
Sol öncül önsavı ispatlar. Sağ öncül önsavı varsayarak asıl sonuca
ulaşır. $A$ formülü sonuçta görünmez; ``kesilip atılmıştır''. Kesme,
tüm diğer kurallardan farklıdır: Sonuçta \emph{bulunmayan} bir formül
(önsav) öncüllerde ortaya çıkar. İspat arama açısından bu bir felakettir.
Kesmeyi aşağıdan yukarıya uygulayacak olsak, sonsuz sayıdaki olası
$A$'lardan birini tahmin etmemiz gerekirdi.

\begin{teorem}{Hauptsatz (Gentzen, 1935)}{hauptsatz}
Kesme kuralı elenebilir: Kesme ile türetilebilen her ardışık, kesme
kullanılmadan da türetilebilir.
\end{teorem}

\noindent Elimizdeki araçlarla bu teoremin önermeler mantığı için
ispatı birkaç satırdır. Kesmenin iki öncülü geçerliyse sonucu da
geçerlidir. Kesmesiz tamlık ise geçerli her ardışığın kesmesiz bir
türetimi olduğunu söyler:

\leandosya[Kesme elenebilir]{Bolum13}{kesme}

\noindent Bu ispat, kesmeyi sisteme bir kural olarak eklememize gerek
olmadığını gösterir. \lc{kesme} bir \emph{teorem} olarak elimizdedir ve
kesmesiz sistemin türetebildiklerini değiştirmez. Böyle kurallara
\emph{kabul edilebilir} (\emph{admissible}) kurallar denir.

\begin{figure}[htbp]
\centering
\begin{tikzpicture}[
    kutu/.style={draw=#1, fill=#1!8, very thick, rounded corners=5pt,
                 text width=2.3cm, align=center, inner sep=6pt, font=\small},
    ok/.style={-{Stealth[length=3mm]}, very thick, #1},
  ]
  \node[kutu=anarenk]     (k) at (0,0)    {Kesmeli türetim\\$\Gamma \Rightarrow \Delta$};
  \node[kutu=ikincil]     (g) at (5.3,0)  {$\Gamma \Rightarrow \Delta$\\geçerli};
  \node[kutu=teoremframe] (s) at (10.6,0) {Kesmesiz türetim\\$\Gamma \Rightarrow \Delta$};
  \draw[ok=ikincil] (k) -- node[above, font=\baslikfont\scriptsize\bfseries]{geçerlilik} (g);
  \draw[ok=teoremframe] (g) -- node[above, font=\baslikfont\scriptsize\bfseries]{kesmesiz tamlık} (s);
  \draw[ok=black!30, dashed] (k.south) to[bend right=18]
       node[below, font=\baslikfont\scriptsize, text=black!60]{Gentzen: türetimi doğrudan dönüştür (sözdizimsel)} (s.south);
\end{tikzpicture}
\caption{Kesme elemenin iki yolu. Üstteki anlambilimsel yol bu bölümde
\Lean{}'de yapıldı. Alttaki sözdizimsel yol, Gentzen'in yöntemidir.}
\label{fig:kesme-yollari}
\end{figure}

\subsection{Sözdizimsel kesme eleme}

Gentzen'in ispatı anlambilim kullanmaz. Kesmeli bir türetimi adım adım
kesmesiz bir türetime \emph{dönüştürür}. Ana fikir, bir kesmeyi
öncüllerine doğru ``yukarı itmektir''. Kesme formülü her iki öncülde de
yeni tanıtılmışsa, kesme daha küçük formüller üzerinde kesmelere bölünür.
Örneğin $A \ove B$ için:
\begin{center}
\scriptsize
\begin{prooftree}
  \hypo{\Gamma \Rightarrow A, \Delta}
  \hypo{\Gamma \Rightarrow B, \Delta}
  \infer2[$\wedge$R]{\Gamma \Rightarrow A \ove B, \Delta}
  \hypo{A, B, \Gamma \Rightarrow \Delta}
  \infer1[$\wedge$L]{A \ove B, \Gamma \Rightarrow \Delta}
  \infer2[Kesme]{\Gamma \Rightarrow \Delta}
\end{prooftree}
\quad\raisebox{1.2em}{\normalsize$\rightsquigarrow$}\quad
\begin{prooftree}
  \hypo{\Gamma \Rightarrow A, \Delta}
  \hypo{\Gamma \Rightarrow B, \Delta}
  \hypo{A, B, \Gamma \Rightarrow \Delta}
  \infer2[Kesme$_B$]{A, \Gamma \Rightarrow \Delta}
  \infer2[Kesme$_A$]{\Gamma \Rightarrow \Delta}
\end{prooftree}
\end{center}
(Sağdaki türetimde öncüllere gereken zayıflatmalar gösterilmemiştir.)
Kesme formülü bir öncülde yeni tanıtılmamışsa, kesme o öncülün kuralının
üstüne taşınır. İspat, kesme formülünün boyutu ve türetimin yüksekliği
üzerinde \emph{iç içe} bir tümevarımdır. Bu tümevarımın \Lean{}'deki
formalizasyonu, \cref{bol:proje}'deki projelerden biridir. Pfenning'in
``yapısal kesme eleme'' yöntemi, bu ispatın ispat asistanlarına uygun bir
biçimidir \citep{pfenning2000}.

\begin{dikkat}[Dikkat: Kesme elemenin bedeli]
Kesme eleme ``bedava'' değildir. Kesmesiz türetimler, kesmeli olanlardan
\emph{akıl almaz} derecede uzun olabilir. Bazı formül aileleri için
kesmeli türetimin boyutu $n$ iken, kesmesiz türetimin boyutu $2^{2^{\cdot^{\cdot^{2}}}}$
(yüksekliği $n$ olan bir üs kulesi) mertebesindedir \citep{statman1979}.
Matematikte önsavların neden bu kadar önemli olduğunun biçimsel bir
açıklaması budur. Önsavlar türetimleri küçültür. Bu bölümün
\lc{ardisik\_tam} ispatı da, kesmesiz türetimleri doğruluk tablosu
büyüklüğünde üretir.
\end{dikkat}

% ---------------------------------------------------------------------
\section{Alt formül özelliği ve sonuçları}\label{sec:alt-formul}

\Cref{fig:g3-kurallar}'deki kurallara bir kez daha bakın. Her kuralın
öncüllerinde geçen her formül, sonucundaki bir formülün \emph{alt
formülüdür}. Kesme bu özelliği taşımayan tek kuraldı.

\begin{teorem}{Alt formül özelliği}{alt-formul}
Kesmesiz bir türetimde geçen her formül, kök ardışıktaki bir formülün alt
formülüdür.
\end{teorem}

\noindent Bu sözdizimsel özelliğin şaşırtıcı derecede güçlü sonuçları
vardır:

\begin{itemize}
  \item \textbf{Tutarlılık.} Boş ardışığın kesmesiz bir türetimi
        olsaydı, alt formül özelliğine göre bu türetimde hiçbir formül
        geçemezdi. Ama Aks ve $\bot$L kurallarının ikisi de bir formül
        gerektirir. Gentzen'in Peano aritmetiği için tutarlılık ispatının
        iskeleti budur. Bölüm~5'te doğal çıkarım için aynı sonuca
        normalleştirmeyle ulaşmıştık.
  \item \textbf{Karar verilebilirlik.} Bir ardışığın olası kesmesiz
        türetimlerinde yalnızca sonlu sayıda formül kullanılabilir. Bu
        yüzden arama sonludur. Önermeler mantığı için bunu
        \lc{tam\_aux}'ta gördük. LJ için de geçerlidir, ama daraltma
        kuralı yüzünden çok daha incedir (aşağıdaki araştırma notuna
        bakın).
  \item \textbf{Ara değer (interpolasyon).} $A \Rightarrow B$ geçerliyse,
        yalnızca $A$ ile $B$'nin \emph{ortak} atomlarını içeren bir $C$
        formülü vardır ki $A \Rightarrow C$ ve $C \Rightarrow B$ geçerlidir
        \citep{craig1957}. Kesmesiz türetim üzerinde tümevarımla $C$
        doğrudan inşa edilebilir.
\end{itemize}

\subsection{Çalıştırılabilir bir ispatlayıcı}

\lc{tam\_aux} bir \emph{teoremdir}; bir türetim olduğunu söyler. Aynı
özyinelemeyi \lc{Bool} döndüren bir \emph{fonksiyon} olarak yazarsak,
çalıştırabileceğimiz bir karar yöntemi elde ederiz:

\leandosya[Çalıştırılabilir ispat arama]{Bolum13}{ara}

\noindent \lc{ara}'nın doğruluğunu (yani \lc{ara} \lc{true} döndürdüğünde
gerçekten bir türetim olduğunu) ispatlamak, \lc{tam\_aux}'un yapısını
izleyen güzel bir alıştırmadır. \ciltiki{6}'da bu fikri \emph{yansıma}
(\emph{reflection}) adıyla genelleştireceğiz. Orada bir karar yönteminin
doğruluğunu bir kez ispatlayıp, sonra onu \leantac{decide} ile çalışan bir
taktik olarak kullanacağız.

% ---------------------------------------------------------------------
\section{Doğal çıkarım ile ardışık hesap}\label{sec:dogal-ardisik}

İki sistemimiz de iki değerli anlambilime göre geçerli ve tamdır. Öyleyse
aynı şeyleri türetirler. \Lean{}'de bu, iki tamlık teoreminin
birleştirilmesiyle tek satırda ispatlanır:

\leandosya[Doğal çıkarım ile denklik]{Bolum13}{denklik}

\noindent Ama iki sistemin ilişkisi ``aynı şeyleri türetmek''ten daha
derindir. \Cref{fig:nd-lk}'de görüldüğü gibi, kurallar arasında
yapısal bir eşleşme vardır.

\begin{figure}[htbp]
\centering
\begin{tikzpicture}[
    kutu/.style={draw=#1, fill=#1!8, thick, rounded corners=4pt,
                 text width=3.6cm, align=center, inner sep=5pt, font=\small,
                 minimum height=1.0cm},
    ok/.style={{Stealth[length=2.4mm]}-{Stealth[length=2.4mm]}, thick, black!40},
  ]
  \node[font=\baslikfont\bfseries, text=anarenk] at (0,1.2) {Doğal çıkarım};
  \node[font=\baslikfont\bfseries, text=ikincil] at (7,1.2) {Ardışık hesap};
  \node[kutu=anarenk] (g)  at (0,0)    {giriş kuralları};
  \node[kutu=anarenk] (e)  at (0,-1.4) {eleme kuralları};
  \node[kutu=anarenk] (d)  at (0,-2.8) {dolambaç\\{\scriptsize (giriş + eleme)}};
  \node[kutu=anarenk] (n)  at (0,-4.2) {normalleştirme};
  \node[kutu=ikincil] (r)  at (7,0)    {sağ kurallar (R)};
  \node[kutu=ikincil] (l)  at (7,-1.4) {sol kurallar (L)};
  \node[kutu=ikincil] (k)  at (7,-2.8) {kesme};
  \node[kutu=ikincil] (ke) at (7,-4.2) {kesme eleme};
  \draw[ok] (g) -- (r); \draw[ok] (e) -- (l);
  \draw[ok] (d) -- (k); \draw[ok] (n) -- (ke);
\end{tikzpicture}
\caption{Doğal çıkarım ile ardışık hesap arasındaki eşleşme. Doğal
çıkarımda bir girişin hemen ardından gelen eleme (dolambaç, Bölüm~5),
ardışık hesapta bir kesmeye karşılık gelir. Prawitz'in normalleştirme
teoremi ile Gentzen'in Hauptsatz'ı aynı madalyonun iki yüzüdür.}
\label{fig:nd-lk}
\end{figure}

\noindent Curry--Howard açısından bakınca bu eşleşme daha da
anlamlıdır. Doğal çıkarım türetimleri $\lambda$-terimleridir;
normalleştirme ise bir programın \emph{çalıştırılmasıdır}. Ardışık hesaptaki
kesme eleme de bir hesaplamadır. Bu hesaplamanın programlama dillerindeki
karşılıkları (örneğin ``\emph{let}'' ifadelerinin açılması) Bölüm~14'te
$\lambda$-hesabı ile birlikte ele alınacak.

% ---------------------------------------------------------------------
\begin{arastirmanotu}
\textbf{Sezgici mantıkta ispat arama.} LJ'de $\to$L kuralı ana formülü
öncüllerde de tutmak zorundadır. Aksi hâlde sistem tam olmaz. Bu yüzden
saf geriye arama sonsuza kadar sürebilir. Roy Dyckhoff'un \emph{LJT}
sistemi \citep{dyckhoff1992}, $\to$L kuralını öncülün biçimine göre dört
kurala bölerek bu sorunu çözer ve her türetimde ölçünün azaldığı,
daraltmasız bir sezgici ardışık hesap verir. Mathlib'deki \leantac{itauto}
gibi sezgici karar yöntemleri bu çizgidedir.

\textbf{Formalize edilmiş kesme eleme.} Kesme eleme, ispat asistanlarında
metamantık formalizasyonunun ölçütlerinden biridir. Chaudhuri, Lima ve
Reis doğrusal mantıklar için kesme elemeyi formalize ettiler
\citep{chaudhuri2019}. Shillito, van~der~Giessen, Goré ve Iemhoff ise
sezgici bir kanıtlanabilirlik mantığı için yeni bir ardışık hesap
önerip kesme elemesini ve ispat aramasının sonlanmasını Coq'ta
formalize ettiler \citep{shillito2023}. Lean~4 tarafında Foundation
kütüphanesi \citep{foundation}, kip mantıkları ve kanıtlanabilirlik
mantığı için çeşitli ispat sistemlerini formalize eder.

\textbf{İspat aramadan otomasyona.} Mathlib'deki \leantac{tauto} ve
genel amaçlı \leantac{aesop} taktiği \citep{limperg2023aesop}, bu
bölümdeki fikirlerin \Lean{}'in kendi mantığına uygulanmış hâlleridir:
Tersinir kurallar ``güvenli'' kabul edilip hemen uygulanır, tersinir
olmayanlar için ise en iyi öncelikli bir arama yapılır.
\end{arastirmanotu}

\begin{bolumozeti}
\begin{itemize}
  \item Bir ardışık $\Gamma \Rightarrow \Delta$, ``$\Gamma$'nın hepsi
        doğruysa $\Delta$'dan biri doğrudur'' demektir. Ardışık hesapta
        her bağlacın bir sol ve bir sağ kuralı vardır.
  \item LK (klasik) ile LJ (sezgici) arasındaki fark, sağ tarafta kaç
        formül bulunabileceğidir. Peirce yasası ve üçüncü hâl, sağda iki
        formül gerektirir.
  \item Sistemimizin her kuralı tersinirdir. Bu yüzden geriye doğru ispat
        arama hiç yanlış seçim yapmaz ve bir ölçü sayesinde sona erer.
  \item Arama, bir tamlık ispatına dönüşür: Geçerli her ardışığın
        kesmesiz bir türetimi vardır (\lc{ardisik\_tam}).
  \item Kesme kuralı elenebilir (\emph{Hauptsatz}). Anlambilimsel ispat,
        geçerlilik ile kesmesiz tamlığın birleşimidir. Gentzen'in
        sözdizimsel ispatı ise türetimleri doğrudan dönüştürür; bunun
        bedeli türetim boyutunda büyük bir artıştır.
  \item Kesmesiz türetimlerin alt formül özelliği; tutarlılık, karar
        verilebilirlik ve interpolasyon sonuçlarının kaynağıdır.
  \item Doğal çıkarım ve ardışık hesap aynı formülleri türetir. Dolambaç
        ile kesme ve normalleştirme ile kesme eleme birbirine karşılık
        gelir.
\end{itemize}
\end{bolumozeti}

\begin{alistirmalar}
\alistirma{1} $p \ove q \Rightarrow q \ove p$. Türetim ağacını da çizin
ve \lc{a11\_3} ile karşılaştırın: Hangisi daha kısa? (\lc{a13\_1})

\alistirma{1} ${\sim}p, p \Rightarrow$\,. Sağ tarafın boş olmasının
anlamı nedir? (\lc{a13\_2})

\alistirma{1} Zayıflatma, yapısal kuralın özel bir hâlidir.
(\lc{a13\_3})

\alistirma{2} Dummett formülü: $\Rightarrow (p \oise q) \oveya (q \oise
p)$. Ana formülü öne almak için \lc{degis}'i kullanmanız gerekecek. Bu
formül sezgici mantıkta türetilemez. Türetiminizin hangi adımı LJ'de
yasaktır? (\lc{a13\_4})

\alistirma{2} $p \oveya q \not\Rightarrow p$: Bu ardışık türetilemez.
(\lc{a13\_5})

\alistirma{2} Sözdizimsel tersinirlik: $A \ove B, \Gamma \Rightarrow
\Delta$ türetilebiliyorsa $A, B, \Gamma \Rightarrow \Delta$ da
türetilebilir. Tamlığı kullanın. (\lc{a13\_6})

\alistirma{2} Doğal çıkarımda ispatlanan bir önsavı, ardışık hesapta
kesme ile kullanın. (\lc{a13\_7})

\alistirma{3} $\Rightarrow A$ türetilebilir $\iff$ $A$ bir totolojidir.
(\lc{a13\_8})

\alistirma{2} (Kâğıt-kalem) $(p \oise q) \oise ({\sim}q \oise {\sim}p)$
için \cref{fig:ispat-arama} benzeri bir arama ağacı çizin. Kaç yaprak
var?

\alistirma{3} (Proje) \lc{ara n Γ Δ P Q = true} ise \lc{Ardisik (Γ ++
P.map atom) (Δ ++ Q.map atom)} olduğunu ispatlayın. İpucu: \lc{tam\_aux}
ile aynı örüntüleri kullanın; atomik durumda \lc{List.any} için bir
önsav gerekecek.

\alistirma{3} (Proje) LJ'yi \lc{ArdisikJ : List Formul → Option Formul
→ Prop} olarak tanımlayın (sağda en fazla bir formül). Geçerliliğini
ispatlayın ve \lc{ArdisikJ [] (some peirce)}'in türetilemediğini
tartışın. Bu türetilemezliği ispatlamak için hangi anlambilim gerekir?
(İpucu: \ciltiki{2}.)

\alistirma{3} (Araştırma) Gentzen'in sözdizimsel kesme eleme ispatını
\lc{Ardisik} sistemine kesme kuralı eklenmiş bir \lc{ArdisikK} için
formalize etmeye çalışın. Tümevarımın ölçüsü ne olmalı? \lc{degis}
kuralı ispatı nasıl etkiliyor?
\end{alistirmalar}
